Two lines meet at $C$ on the vertical axis, so their vertical intercepts coincide. Since the fixed intercepts six and twelve differ, $b$ is six or twelve. Also $a\ne0$, since two horizontal nonzero-intercept lines could not provide two lines through $B$ on the horizontal axis.
If $b=6$, the lines through $C$ are the first two. The first and third have different horizontal intercepts, so $B$ is the intersection of the second and third on the horizontal axis. Their intercepts give $-2=-\dfrac{24}{a}$, so $a=12$. The remaining intersection is $A=(1,18)$, with sum 19.
If $b=12$, the second and third meet at $C$. Equating horizontal intercepts of the first and second gives $-\dfrac{6}{a}=-4$, so $a=\dfrac{3}{2}$. The remaining intersection is $A=(8,18)$, with sum 26. Both are valid nondegenerate triangles, so the minimum is $\boxed{19}$.
\begin{asy}
texpreamble("\usepackage{amsmath,amssymb}");
size(240);
defaultpen(fontsize(9pt));
pair A=(1,18),B=(-2,0),C=(0,6);draw((-2.5,0)--(2,0),Arrow);draw((0,-1)--(0,19),Arrow);draw(A--B--C--cycle,blue);label("$A(1,18)$",A,(0,1));label("$B(-2,0)$",B,(-1,-1));label("$C(0,6)$",C,(1,0));
\end{asy}
